\section{A self test for the Pauli basis}\label{sec:self-test-pauli}

In this section, we give a self test for the Pauli basis measurement.
Given~$W \in \{X, Z\}$, this test compels the prover
to measure an EPR register in the~$W$ basis
and return the outcome to the verifier.

\begin{definition}
The \emph{Pauli basis strategy with parameters $n$ and $q$} (a prime power),
denoted $\paulistrat{n}{q}$, is the partial strategy
with the state $\ket{\mathrm{EPR}_q^n}$ and measurement matrices $\tau^W_u$ for each $W \in \{X,Z\}, u \in \F_q^n$.
\end{definition}

The main result of this section is the following self-test for the case when~$q$ is a power of~$2$.

\begin{theorem}\label{thm:basis-test}
Let $\bW \sim \{X, Z\}$ uniformly at random.  Let $n$, $q$, $\eta$ satisfy the Pauli basis condition.
Then there is a self-test $\game_{\mathrm{basis}} := \game_{\mathrm{basis}}(n,q)$ for $\paulistrat{n}{q}$ over~$\bW$
with robustness $\delta(\eps) = \poly(\eps, \eta)$.
Moreover, there is a value-$1$ real commuting EPR strategy with auxiliary state $\ket{\mathrm{EPR}_2}$.
Finally,
\begin{equation*}
\qlength{\game_{\mathrm{basis}}} = O\left(\log(n)\right),
\quad
\alength{\game_{\mathrm{basis}}} = \poly(n),
\end{equation*}
\begin{equation*}
\qtime{\game_{\mathrm{basis}}} = O\left(\log(n)\right),
\quad
\atime{\game_{\mathrm{basis}}} = \poly(n).
\end{equation*}
\end{theorem}

We prove this by a straightforward reduction to the quantum low-degree test of~\cite{NV18a}.

\subsection{The quantum low-degree test}

The goal of the quantum low-degree test of~\cite{NV18a}
is to force the provers to use a ``compressed" version of the Pauli basis strategy.
Given~$\bW$, they should measure their register in the $\bW$ basis, receiving $\bu \in \F_q^n$.
However, $\bu$, a length-$n$ string, might be prohibitively expensive to communicate to the verifier,
so they should instead compute the low degree encoding $g_{\bu}$
and return its evaluation at a single point $\bw \in \F_q^m$ of the verifier's choosing.
(The point of this section is to ``uncompress" their protocol.)

\begin{definition}
Fix parameters for the low-degree encoding $\params := (q = p^t, h, H, m, n, \pi)$
satisfying the ``low-degree conditions"
$h \leq q$, and $n \leq h^m$.
For any string $u \in \F_q^n$,
these parameters give a low-degree encoding $g_u:\F_q^m \rightarrow \F_q$.

The \emph{low-degree Pauli strategy with parameters~$\params$},
denoted $\mathcal{LD}(\params)$, is the partial strategy
with state $\ket{\mathrm{EPR}^n_q}$ and measurement matrices
\begin{equation*}
\tau^{W, w}_a :=
\tau^{W}_{[g_u(w) = a]} =
\sum_{u : g_u(w) = a} \tau^W_u
\end{equation*}
for each $W \in \{X, Z\}, w \in \F_q^m, a \in \F_q$.
Equivalently, this is the strategy where we perform the Pauli $W$-basis measurement
and output the low-degree encoding of the outcome~$u$ evaluated at the point~$w$, i.e.\ the value $g_u(w)$.
\end{definition}

The main result of~\cite{NV18a} is the following.

\begin{theorem}[{\cite[Theorem 3.2]{NV18a}}]\label{thm:anand-thomas-low-degree}
Fix low-degree parameters~$\params$ with $p = 2$ (so that $q = 2^t$) and $m \geq 2$, 
and let $\calD$ be the uniform distribution over $(W, w)$ with $W \in \{X, Z\}, w \in \F_q^m$.
Then there is a self-test $\game_{\mathrm{Qlowdeg}} := \game_{\mathrm{Qlowdeg}}(\params)$
for $\mathcal{LD}(\params)$ over~$\calD$
with robustness
$
\delta(\epsilon) = \mathrm{poly}(\epsilon,md/q^c),
$
with $c > 0$.
%\ainnote{TODO: Propagate to the rest of the document!!}
Moreover, there is a value-$1$ real commuting EPR strategy with auxiliary state $\ket{\mathrm{EPR}_2}$.
Finally,
\begin{equation*}
\qlength{\game_{\mathrm{Qlowdeg}}} = O(m \log q),
\quad
\alength{\game_{\mathrm{Qlowdeg}}} = O(d^2 \log (q)),
\end{equation*}
\begin{equation*}
\qtime{\game_{\mathrm{Qlowdeg}}} = O(m \log q),
\quad
\atime{\game_{\mathrm{Qlowdeg}}} = \poly(m, d, \log q).
\end{equation*}
\end{theorem}

(We note that this result is stated in \cite{NV18a} for general primes~$p$.
However, the $p \neq 2$ case relied on a self-testing result for a generalization of the Magic Square game
which was recently discovered to contain a bug. 
Fortunately, the $p=2$ case needs only a self-testing result for the
``traditional" binary Magic Square game, and this follows
from~\cite{WBMS16}.)

\begin{remark}
We note that the quantum low-degree test, as stated in \cite{NV18a}, does \emph{not} have value-$1$ real commuting EPR strategies.
This is because it uses as a subroutine the standard magic square game,
and the magic square game does not have value-$1$ real commuting EPR strategies.
Its value-$1$ strategies \emph{are} EPR strategies,
and they \emph{are} real (all observables are either $X$ or $Z$, with the sole exception of the $Y \otimes Y$ observable,
which can be rewritten as $Y \otimes Y = -(X \otimes X) \cdot (Z \otimes Z)$, manifestly real).
But they are not commuting, because each row and column have at least one pair of noncommuting observables.

Consider instead the following ``oracularized" version of the magic square game:
one player is given a random row or column (and is expected to play as in the normal magic square game),
and the other player is given a random cell in that row or column,
and the verifier simply checks that they agree on that cell.
In addition, with some constant probability, both players are given the same cell and their answers are checked against each other.
In this case, all observables measured are commuting, and so this variant has a value-$1$ real commuting EPR strategy.
In addition, it certifies the same state and measurements as the normal magic square game,
and so we can use it as a subroutine in the quantum low-degree test instead.
\end{remark}

\begin{remark}
We note again that the soundness case in our definition of a self-test is quite different from the one given in \cite[Definition~$2.5$]{NV18a},
and it is not clear that a self-test in their sense implies a self-test in our sense.
However, for the quantum low-degree test, their soundness case \emph{does} match ours.
By~\cite[Lemma~$4.1$]{NV18a}, there is a local isometry~$\phi = \phi_1 \otimes \phi_2$ such that
\begin{equation}\label{eq:anand-testing-state}
\Vert \phi \ket{\overline{\psi}} - \ket{\psi}\ket{\mathrm{aux}} \Vert^2 \leq \delta(\epsilon).
\end{equation}
and
\begin{equation}\label{eq:anand-testing-state}
\E_{(\bW, \bw)} \sum_a \Vert \phi \cdot (\overline{M}^{\bW, \bw}_{a} \otimes I_{\reg{Bob}}) \ket{\overline{\psi}}
	 - (\tau_{a}^{\bW, \bw} \otimes I_{\reg{aux}})  \otimes I_{\reg{Bob}} \ket{\psi} \ket{\mathrm{aux}} \Vert^2
\leq \delta(\eps).
\end{equation}
The key difference from our self-test definition is that, as stated, their local isometry need not be symmetric (i.e.\ $\phi_1 \neq \phi_2$),
but their construction actually \emph{does} give a symmetric isometry with $\phi_1 = \phi_2$.
Then, from \Cref{eq:anand-testing-state} it is easy to derive \Cref{eq:cite-once-then-never-again} using \Cref{eq:anand-testing-state} and the triangle inequality (\Cref{fact:triangle}).
\end{remark}

\subsection{Proof of \Cref{thm:basis-test}: the Pauli basis test}

We now state the Pauli basis test.

\begin{definition}
Let $n,q,\eta$ be as in \Cref{thm:basis-test}.
Fix the remaining low-degree parameters~$\params$ as follows:
\begin{equation*}
h = \lceil q^{1/2}\rceil, \qquad m = 2\cdot \left\lceil \frac{\log(n)}{\log(q)}\right\rceil,
\qquad d = m \cdot (h-1).
\end{equation*}
Then the \emph{Pauli basis game}~$\game_{\mathrm{basis}}(n,q)$ is given by~\Cref{fig:pauli-basis}.
\end{definition}

{
\floatstyle{boxed} 
\restylefloat{figure}
\begin{figure}
With probability~$\tfrac{1}{2}$ each, perform one of the following two tests.
\begin{enumerate}
\item \textbf{Low-degree:} Perform $\game_{\mathrm{Qlowdeg}}(\params)$. \label{item:qlowdeg-test}
\item \textbf{Cross-check:}  Draw $\bW \sim \{X, Z\}$, $\bw \sim \F_q^m$. Flip an unbiased coin $\bb \sim \{0, 1\}$.
	Distribute the questions as follows:
	\begin{itemize}
	\item[$\circ$] Player~$\bb$: Give $\bW$; receive $\bu \in \F_q^n$.
	\item[$\circ$] Player~$\overline{\bb}$: give $(\bW,\bw)$; receive $\ba$.
	\end{itemize}
	Accept if $g_{\bu}(\bw) = \ba$.
\end{enumerate} \label{item:check-with-qlowdeg}
\caption{The game $\game_{\mathrm{basis}}(n,q)$.\label{fig:pauli-basis}}
\end{figure}
}

These parameters are chosen so that they are valid low-degree parameters
(guaranteeing the existence of the low-degree code),
which is necessary for the quantum low-degree test.
In particular, these satisfy (i) $h \leq q$ and (ii) $n \leq h^m$.
The first of these is immediate; as for the second,
\begin{equation*}
h^m \geq (q^{1/2})^{2 \cdot \log(n)/\log(q)} = q^{\log(n)/\log(q)} = n.
\end{equation*}
In addition, the code has relative distance $d/q \leq mh/q \leq \eta$.
\begin{equation*}
\frac{d}{q} = \frac{m \cdot(h-1)}{q} \leq \frac{m h}{q} \leq 8\cdot \frac{\log(n)}{\log(q)} \cdot \frac{q^{1/2}}{q} \leq \frac{8 \log(n)}{q^{1/2}} \leq \eta,
\end{equation*}
where the final step is because $n, q, \eta$ satisfy the Pauli basis condition.
Finally, we note that even if~$q$ is a large polynomial of~$n$, $m$ is always at least~$2$, which permits us to use the quantum low-degree test.
We now prove \Cref{thm:basis-test}.

\begin{proof}[Proof of \Cref{thm:basis-test}]
The question lengths and times of both the quantum low-degree test and the cross-check are given by
\begin{equation*}
m \log(q) = 2\cdot \left\lceil \frac{\log(n)}{\log(q)}\right\rceil  \cdot \log(q) = O(\log(n)).
\end{equation*}
As for the answer lengths and times, these are bounded by $\poly(n)$ for both the quantum low-degree test and the cross-check.
We now consider the completeness and soundness cases separately.
\paragraph{Completeness.}
Let~$(\psi, M)$ be the value-$1$ commuting EPR strategy for the quantum low-degree test guaranteed by \Cref{thm:anand-thomas-low-degree}.
This has state $\ket{\psi} = \ket{\mathrm{EPR}_q^n} \ket{\mathrm{EPR}_2}$ and measurement matrices $M^{W,w}_a = \tau^{W,w}_a \otimes I_{\reg{aux}}$.
If we add in the measurement matrices $M^W_u = \tau^W_u \otimes I_{\reg{aux}}$,
then this strategy passes the cross-check with probability~$1$.
This is because after Player~$\bb$ measures~$\bu$,
the state collapses to $\ket{\tau^{\bW}_{\bu}}\ket{\tau^{\bW}_{\bu}} \ket{\mathrm{EPR}_2}$,
and so Player~$\overline{\bb}$ will measure $\ba = g_{\bu}(\bw)$.
As a result, this is a value-$1$ strategy.
Furthermore, it is a commuting EPR strategy because the cross-check measurements $M^W$ and $M^{W,w}$ commute.
Finally, this strategy extends the Pauli basis strategy.
This proves the completeness case.

\paragraph{Soundness.}
Throughout the soundness, we will use $\delta(\eps)$ to denote a function of the form $\poly(\eps, \eta)$ which may change from use to use.
The $\delta(\eps)$ in \Cref{thm:anand-thomas-low-degree} is of this form because $d/q\leq \eta$.

Let $\overline{\calS} = (\overline{\psi}, \overline{M})$ be a strategy with $\valstrat{\game_{\mathrm{basis}}}{\overline{\calS}} = 1-\eps$.
Then this strategy must pass $\game_{\mathrm{Qlowdeg}}$ with probability at least $1-2\eps$.
By \Cref{thm:anand-thomas-low-degree} this gives us a local isometry
$\phi = \phi_{\mathrm{local}}\otimes \phi_{\mathrm{local}}$ and a state $\ket{\mathrm{aux}}$
with the following properties:
if we define the new strategy $\calS$
in which $\ket{\psi} = \phi \ket{\overline{\psi}}$
and $M_a^x = \phi_{\mathrm{local}}\cdot \overline{M}_a^x \cdot \phi_{\mathrm{local}}^\dagger$, then
\begin{equation}\label{eq:cant-think-of-a-good-name}
\Vert \ket{\psi} - \ket{\mathrm{EPR}_q^n} \ket{\reg{aux}} \Vert^2 \leq \delta(\eps),
\qquad
M_a^{W,w} \otimes I_{\reg{Bob}}
\approx_{\delta(\eps)}
(\tau_a^{W,w} \otimes I_{\reg{aux}}) \otimes I_{\reg{Bob}},
\end{equation}
on state $\ket{\psi}$ and distribution $\calD$.
Because $\calS$ is just a rotated version of $\overline{\calS}$,
it also passes $\game_{\mathrm{basis}}$ with probability $1-\eps$.
As a result, $\calS$ passes the cross-check in \Cref{item:check-with-qlowdeg} with probability at least $1-2\eps$.
By \Cref{fact:agreement}, we conclude that
\begin{equation}\label{eq:worth-studying}
M^W_{[g_u(w)=a]} \otimes I_{\reg{Bob}}
\approx_{\eps}
I_{\reg{Alice}} \otimes M^{W,w}_a
\approx_{\delta(\eps)} I_{\reg{Alice}} \otimes (\tau^{W,w}_a \otimes I_{\reg{aux}})
= I_{\reg{Alice}} \otimes (\tau^{W}_{[g_v(w)=a]} \otimes I_{\reg{aux}})
\end{equation}
on state $\ket{\psi}$.
By \Cref{fact:almost-agreement} and the fact that the~$\tau$ measurements are projective, this implies that
\begin{equation*}
M^W_{[g_u(w)=a]} \otimes I_{\reg{Bob}}
\consistency_{\delta(\eps)} I_{\reg{Alice}} \otimes (\tau^{W}_{[g_v(w)=a]} \otimes I_{\reg{aux}})
\end{equation*}
Now by
\Cref{prop:same-on-point-same-on-subspace} (where we let~$\bs$ be the singleton distribution on the ``trivial" subspace $\bs = \F_q^m$)
and the fact that $d/q \leq \eta$, we can conclude that
\begin{equation*}
M^W_{u} \otimes I_{\reg{Bob}}
\consistency_{\delta(\eps)} I_{\reg{Alice}} \otimes (\tau^{W}_{u} \otimes I_{\reg{aux}}).
\end{equation*}
Applying \Cref{fact:agreement} again, this yields
\begin{equation}\label{eq:almost-there-almost-there}
M_{u}^{W} \otimes I_{\reg{Bob}}
\approx_{\delta(\eps)} I_{\reg{Alice}} \otimes (\tau_{u}^{W} \otimes I_{\reg{aux}})
\approx_{\delta(\eps)} (\tau_{u}^W \otimes I_{\reg{aux}}) \otimes I_{\reg{Bob}}
\end{equation}
on state $\ket{\psi}$, where the last step uses \Cref{fact:the-ol-state-y-swaperoonie}
to combine \Cref{fact:the-ol-pauli-swaperoonie} with \Cref{eq:cant-think-of-a-good-name}.
The analogous statement for the state $\ket{\mathrm{EPR}_q^n}$ follows from~\Cref{fact:the-ol-state-y-swaperoonie}.
This establishes the theorem.
\end{proof}

  
%%% Local Variables:
%%% mode: latex
%%% TeX-master: "main"
%%% End:
%%%---------- close: pauli-basis

%%%%%%%%%%% jump to pauli-basis-layer
%%%---------- open: pauli-basis-layer
%!TEX root = main.tex

